\section{Descomposición en \(5\) modos de polarización}

La dimensi\'on cinco adquiere contenido tensorial al elegir una direcci\'on
unitaria de propagaci\'on \(n\) y un marco ortonormal orientado
\((e_1,e_2,n)\). Definimos los cinco tensores reales
\begin{align}
E_{+}&=\frac1{\sqrt2}(e_1e_1^T-e_2e_2^T),
&E_{\times}&=\frac1{\sqrt2}(e_1e_2^T+e_2e_1^T),
\label{rm:eq:gravity-helicity2-real}\\
E_{1}&=\frac1{\sqrt2}(e_1n^T+ne_1^T),
&E_{2}&=\frac1{\sqrt2}(e_2n^T+ne_2^T),
\label{rm:eq:gravity-helicity1-real}\\
E_{0}&=\frac1{\sqrt6}(2nn^T-e_1e_1^T-e_2e_2^T).
\label{rm:eq:gravity-helicity0-real}
\end{align}
Todos son sim\'etricos y sin traza. Los dos primeros son transversales; los
dos siguientes mezclan el plano transversal con la direcci\'on de
propagaci\'on, y el \`ultimo es el modo escalar longitudinal sin traza.

\begin{theorem}[Descomposici\'on helicoidal completa del portador]
\label{rm:thm:gravity-five-helicities}
Los tensores de
\cref{rm:eq:gravity-helicity2-real,rm:eq:gravity-helicity1-real,rm:eq:gravity-helicity0-real}
forman una base ortonormal de \(\operatorname{STF}_3\). En la
complejificaci\'on, las combinaciones
\begin{equation}
\label{rm:eq:gravity-complex-helicities}
E_{\pm2}=\frac1{\sqrt2}(E_+\mp\ii E_\times),
\qquad
E_{\pm1}=\frac1{\sqrt2}(E_1\mp\ii E_2),
\qquad E_0=E_0
\end{equation}
son vectores de peso \(\pm2,\pm1,0\), respectivamente, bajo rotaciones
alrededor de \(n\). Adem\'as,
\begin{equation}
\label{rm:eq:gravity-tt-on-five}
\Lambda(n)E_{\pm2}=E_{\pm2},
\qquad
\Lambda(n)E_{\pm1}=0,
\qquad
\Lambda(n)E_0=0.
\end{equation}
\end{theorem}

\begin{proof}
En el marco \((e_1,e_2,n)\), cada tensor tiene norma de Frobenius uno.
Sus patrones de entradas no diagonales son disjuntos, y las tres diagonales
de \(E_+,E_0\) tienen producto interno cero; por tanto, los cinco tensores
son ortonormales. Como \(\dim\operatorname{STF}_3=6-1=5\), forman una base.

Una rotaci\'on de \`angulo \(\psi\) alrededor de \(n\) env\'ia
\((e_1,e_2)\) a
\((\cos\psi\,e_1+\sin\psi\,e_2,
-\sin\psi\,e_1+\cos\psi\,e_2)\). Sustituyendo en las f\'ormulas se obtiene
una rotaci\'on de \`angulo \(2\psi\) sobre
\(\operatorname{span}\{E_+,E_\times\}\), una de \`angulo \(\psi\) sobre
\(\operatorname{span}\{E_1,E_2\}\), y \(E_0\) fijo. De ah\'i los pesos de
\cref{rm:eq:gravity-complex-helicities}. Finalmente,
\(\Pi E_{+}\Pi=E_+\) y \(\Pi E_\times\Pi=E_\times\), mientras
\(\Pi E_1\Pi=\Pi E_2\Pi=0\). Para \(E_0\),
\(\Pi E_0\Pi=-(e_1e_1^T+e_2e_2^T)/\sqrt6\), que es puramente traza en
el plano, y la sustracci\'on de \cref{rm:eq:gravity-lambda-tt} la anula.
\end{proof}

El teorema precisa el significado de ``cinco polarizaciones'' sin atribuir
a todas el mismo estatuto din\'amico. Antes de imponer ecuaciones de campo,
el portador de esp\'in dos tiene cinco coordenadas irreducibles. En una
realizaci\'on masiva de Fierz--Pauli, las restricciones dejan precisamente
los cinco pesos de \cref{rm:eq:gravity-complex-helicities}; en una realizaci\'on
sin masa con invariancia de calibre, los sectores \(\pm1\) y \(0\) se eliminan y
quedan \(\pm2\). HMT proporciona aqu\'i el portador, la base y el proyector;
la elecci\'on de la din\'amica masiva o sin masa pertenece a la interfaz
f\'isica. Por eso cinco no significa ``cinco ondas observadas'', sino cinco
grados tensoriales disponibles antes de la reducci\'on.

La orientaci\'on tambi\'en queda visible: cambiar
\((e_1,e_2,n)\) por \((e_1,-e_2,-n)\) conserva \(E_+,E_0\) e intercambia
las fases de las parejas helicoidales. Esa informaci\'on no est\'a en el
escalar \(G_a\); reside en el portador sobre el que act\'ua
\(\mathbb G_{5,a}\).

\section{Cadena representacional \(12\to5\to5\to2\): portador \(V_5\), tensor sin traza, polarizaciones de Fierz–Pauli y helicidades transversales}

La dimensi\'on cinco aparece en objetos distintos, y s\'olo
\cref{rm:eq:gravity-12-5-5-2} autoriza el transporte entre ellos:

\begin{enumerate}
\item \(V_5\) es el sumando irreducible de la representaci\'on sobre los
doce v\'ertices. Es un resultado finito de teor\'ia de representaciones.
\item \(\operatorname{STF}_3\) es el portador real de dimensi\'on cinco de
los pesos \(m=-2,-1,0,1,2\) de esp\'in dos.
\item Una teor\'ia masiva de Fierz--Pauli puede realizar esos cinco pesos
como cinco estados f\'isicos, despu\'es de fijar acci\'on y restricciones.
\item La relatividad general sin masa conserva dos helicidades radiativas,
\(\pm2\), representadas por \(E_+\) y \(E_\times\).
\item Las clases \(E(2)\) de ondas m\'etricas, como \(III_5\), clasifican
posibles amplitudes bajo otras hip\'otesis; no son autom\'aticamente
\(V_5\), Fierz--Pauli o relatividad general.
\end{enumerate}

En particular, la frase correcta es: \emph{HMT construye un portador
tensorial de cinco componentes antes de las restricciones; la reducci\'on
TT de la realizaci\'on sin masa posee rango dos}. No se afirma que la
relatividad general tenga cinco polarizaciones propagantes.
